
\documentclass{article}
\usepackage{colortbl}
\usepackage{makecell}
\usepackage{multirow}
\usepackage{supertabular}

\begin{document}

\newcounter{utterance}

\centering \large Interaction Transcript for game `cladder', experiment `full\_v1.5\_default', episode 3598 with Qwen3.5{-}9B.
\vspace{24pt}

{ \footnotesize  \setcounter{utterance}{1}
\setlength{\tabcolsep}{0pt}
\begin{supertabular}{c@{$\;$}|p{.15\linewidth}@{}p{.15\linewidth}p{.15\linewidth}p{.15\linewidth}p{.15\linewidth}p{.15\linewidth}}
    \# & \multicolumn{2}{c}{Player} && \multicolumn{2}{c}{Game Master} \\
    \hline

    \theutterance \stepcounter{utterance}  
    & & & \multicolumn{4}{p{0.6\linewidth}}{
        \cellcolor[rgb]{0.9,0.9,0.9}{
            \makecell[{{p{\linewidth}}}]{
                \texttt{\tiny{[P1$\langle$GM]}}
                \texttt{You are an expert in causal inference. The following question is not a typical commonsense query, but rather a meticulously designed question created by a professor specializing in causal inference, intended to assess the students' mastery of the course content.} \\
\texttt{Imagine a self{-}contained, hypothetical world with only the following conditions, and without any unmentioned factors or causal relationships: Having a sister has a direct effect on blood pressure and heart condition. Blood pressure has a direct effect on heart condition. For people who do not have a sister and with low blood pressure, the probability of healthy heart is 63\%. For people who do not have a sister and with high blood pressure, the probability of healthy heart is 31\%. For people who have a sister and with low blood pressure, the probability of healthy heart is 72\%. For people who have a sister and with high blood pressure, the probability of healthy heart is 41\%. For people who do not have a sister, the probability of high blood pressure is 66\%. For people who have a sister, the probability of high blood pressure is 38\%. Does having a sister positively affect heart condition through blood pressure?} \\
\texttt{Start your answer with 'Yes' or 'No', followed by additional reasoning or evidence to support your explanation.} \\
            }
        }
    }
     \\ \\

    \theutterance \stepcounter{utterance}  
    & \multicolumn{4}{p{0.6\linewidth}}{
        \cellcolor[rgb]{0.561,0.737,0.561}{
            \makecell[{{p{\linewidth}}}]{
                \texttt{\tiny{[P1$\rangle$GM]}}
                \texttt{No} \\
\\ 
\texttt{**Reasoning:**} \\
\\ 
\texttt{To determine if having a sister positively affects heart condition *through* blood pressure, we must calculate the **Natural Direct Effect (NDE)** and the **Natural Indirect Effect (NIE)** of having a sister on heart condition, mediated by blood pressure.} \\
\\ 
\texttt{Let:} \\
\texttt{*   \$S\$ = Sister (1 = has sister, 0 = no sister)} \\
\texttt{*   \$B\$ = Blood Pressure (1 = high, 0 = low)} \\
\texttt{*   \$H\$ = Heart Condition (1 = healthy, 0 = unhealthy)} \\
\\ 
\texttt{**1. Establish the Probabilities:**} \\
\texttt{*   \$P(B=1 | S=0) = 0.66 \textbackslash{}implies P(B=0 | S=0) = 0.34\$} \\
\texttt{*   \$P(B=1 | S=1) = 0.38 \textbackslash{}implies P(B=0 | S=1) = 0.62\$} \\
\\ 
\texttt{*   \$P(H=1 | S=0, B=0) = 0.63\$} \\
\texttt{*   \$P(H=1 | S=0, B=1) = 0.31\$} \\
\texttt{*   \$P(H=1 | S=1, B=0) = 0.72\$} \\
\texttt{*   \$P(H=1 | S=1, B=1) = 0.41\$} \\
\\ 
\texttt{**2. Calculate the Natural Indirect Effect (NIE):**} \\
\texttt{The NIE measures the change in the outcome (\$H\$) when the mediator (\$B\$) changes from the level it would have taken under the counterfactual "no sister" scenario to the level it actually takes under the "has sister" scenario, while holding the treatment (\$S\$) fixed at "has sister" (to isolate the pathway through \$B\$).} \\
\\ 
\texttt{*   **Counterfactual Mediator State (\$B\_\{S=0\}\$):** If the person did not have a sister, the probability of high blood pressure is 0.66. Thus, the probability of *low* blood pressure is \$1 {-} 0.66 = 0.34\$.} \\
\texttt{*   **Actual Mediator State (\$B\_\{S=1\}\$):** If the person has a sister, the probability of high blood pressure is 0.38. Thus, the probability of *low* blood pressure is \$1 {-} 0.38 = 0.62\$.} \\
\\ 
\texttt{We compare the outcome probability for a person with a sister (\$S=1\$) under two scenarios:} \\
\texttt{1.  **Actual Mediator:** Low blood pressure (since \$P(B=0|S=1) = 0.62\$).} \\
\texttt{    \$P(H=1 | S=1, B=0) = 0.72\$} \\
\texttt{2.} \\
            }
        }
    }
    & & \\ \\

    \theutterance \stepcounter{utterance}  
    & & & \multicolumn{2}{p{0.3\linewidth}}{
        \cellcolor[rgb]{0.9,0.9,0.9}{
            \makecell[{{p{\linewidth}}}]{
                \texttt{\tiny{[GM$|$GM]}}
                \texttt{No} \\
            }
        }
    }
    & & \\ \\

    \theutterance \stepcounter{utterance}  
    & & & \multicolumn{2}{p{0.3\linewidth}}{
        \cellcolor[rgb]{0.9,0.9,0.9}{
            \makecell[{{p{\linewidth}}}]{
                \texttt{\tiny{[GM$|$GM]}}
                \texttt{yes} \\
            }
        }
    }
    & & \\ \\

    \theutterance \stepcounter{utterance}  
    & & & \multicolumn{2}{p{0.3\linewidth}}{
        \cellcolor[rgb]{0.9,0.9,0.9}{
            \makecell[{{p{\linewidth}}}]{
                \texttt{\tiny{[GM$|$GM]}}
                \texttt{game\_result = LOSE} \\
            }
        }
    }
    & & \\ \\

\end{supertabular}
}

\end{document}
